\documentclass[11pt]{article}
\usepackage{mathunal-base}
\mathunalfuente{serif}
\newlist{opciones}{enumerate}{1}
\setlist[opciones]{label=\alph*.,itemsep=1pt,topsep=2pt,leftmargin=1.8em,font=\normalfont}
\newlist{romanitos}{enumerate}{1}
\setlist[romanitos]{label=\roman*.,itemsep=1pt,topsep=2pt,leftmargin=1.8em,font=\normalfont}

\renewcommand{\examtitulo}{3002192 Métodos Numéricos --- Primer Examen Parcial}
\renewcommand{\examfecha}{6 de Marzo de 2006}

\begin{document}

\examheader

\vspace{6pt}
\noindent Nombre \dotfill\ Carné \dotfill\ Profesor \dotfill\ Grupo \dotfill

\vspace{10pt}
\noindent\textbf{1.} Considere la ecuación
\[
x-2+\ln(x)=0, \qquad \text{con } x>0.
\]
\begin{opciones}
\item (5\%) Compruebe que hay un único número real positivo $p$ que es raíz de esta ecuación.
\item (5\%) Verifique que se puede aplicar el método de Bisección a la función
\[
f(x)=x-2+\ln(x)
\]
en el intervalo $[1.5,1.6]$ para aproximar la raíz $p$.
\item (5\%) Calcule las dos primeras iteraciones del método de bisección.
\item (15\%) Verifique que la función $g(x)=2-\ln(x)$ es tal que:
\begin{romanitos}
\item $x=g(x)$ si y solo si $f(x)=0$.
\item $g(x)$ satisface las hipótesis del Teorema de Punto Fijo en $[1.5,1.6]$.
\end{romanitos}
\item (5\%) Calcule las dos primeras iteraciones de punto fijo con 1.55 como aproximación inicial. ¿Cuál es el valor aproximado de la raíz $p$?
\end{opciones}

\vspace{5mm}
\noindent\textbf{2.} Considere la ecuación polinómica
\[
p(x)=x^3+3x-1=0.
\]
Al dividir $p(x)$ por el factor cuadrático $(x^2-ux-v)$ resulta
\[
p(x)=(x^2-ux-v)q(x)+r(x),
\]
con cociente y residuo dados por
\[
q(x)=b_nx^{n-2}+b_{n-1}x^{n-3}+\ldots+b_3x+b_2 \qquad \text{y} \qquad r(x)=b_1(x-u)+b_0
\]
respectivamente.
\begin{opciones}
\item (20\%) Calcule la primera iteración del método de Bairstow para aproximar un factor cuadrático $x^2-ux-v$ del polinomio $p(x)$. Utilice $u_0=-1$ y $v_0=-3$ como aproximaciones iniciales.
\item (10\%) Suponga que el valor de los coeficientes del factor cuadrático es $u=-0.32$ y $v=-3.1$. Encuentre las tres raíces de la ecuación $p(x)=0$.
\end{opciones}

\vspace{5mm}
\noindent\textbf{3.} Considere el sistema lineal
\[
\begin{aligned}
4x_1+2x_2+2x_3&=1\\
x_1+5x_2-x_3&=1\\
2x_1+x_2-3x_3&=1.
\end{aligned}
\]
Se sabe que:
\begin{opciones}
\item[i.] El polinomio característico de la matriz de iteración de Jacobi $T_J$ es $p_{T_J}(x)=-x^3-\dfrac{1}{6}x-\dfrac{1}{30}$.
\item[ii.] La matriz de iteración del método de Gauss-Seidel es
\[
T_G=\begin{bmatrix} 0 & -0.5 & -0.5 \\ 0 & 0.1 & 0.3 \\ 0 & -0.3 & -0.7 \end{bmatrix}.
\]
\item[iii.] El polinomio característico de $T_G$ es $p_{T_G}(x)=-x\left(x^2+\dfrac{2}{15}x+\dfrac{1}{15}\right)$.
\end{opciones}
\begin{opciones}
\item (10\%) Analice la convergencia del método de Jacobi para este sistema.
\item (10\%) Analice la convergencia del método de Gauss-Seidel para este sistema.
\item (7\%) Escriba las fórmulas escalares de iteración del método de Jacobi para el sistema dado y calcule una iteración con aproximación inicial $X^{(0)}$ dada por el vector nulo.
\item (8\%) Escriba las fórmulas escalares de iteración del método de Gauss-Seidel para el sistema dado y calcule una iteración con aproximación inicial $X^{(0)}$ dada por el vector nulo.
\end{opciones}

\vfill
\rule{\linewidth}{0.4pt}\\[1pt]
{\scriptsize\textit{Transcripción digital --- MathUNAL}\hfill\textcolor{unalblue}{\textbf{1}}}

\end{document}
